\documentclass[12pt, a4paper]{article}%Tipo de documento y tamaño de letra%
\usepackage[spanish]{babel}
\usepackage{lipsum}
\usepackage{color,graphicx}%Para poder insertar graficas%
\usepackage[ruled,vlined]{algorithm2e}
\usepackage{amsmath,amssymb}%Insertar unos símbolos matemáticos especiales%
%%%%%%%%%%%%%%%% LA INTEGRAL COMEMIERDA ESA%%

\def\Xint#1{\mathchoice
{\XXint\displaystyle\textstyle{#1}}%
{\XXint\textstyle\scriptstyle{#1}}%
{\XXint\scriptstyle\scriptscriptstyle{#1}}%
{\XXint\scriptscriptstyle\scriptscriptstyle{#1}}%
\!\int}
\def\XXint#1#2#3{{\setbox0=\hbox{$#1{#2#3}{\int}$ }
\vcenter{\hbox{$#2#3$ }}\kern-.6\wd0}}
\def\ddashint{\Xint=}
\def\dashint{\Xint-}

%%%%%%%%%%%%%%%%%%%%%%%%%%%%%%%%%%%%%%%%%%%%
%para poner código xd
\usepackage{listings}
\lstset{
    language=Matlab,
    basicstyle=\ttfamily\small,
    numbers=left,
    stepnumber=1,
    numbersep=10pt,
    frame=single,
    breaklines=true,
    showstringspaces=false,
    keywordstyle=\color{blue!75!black},
    commentstyle=\color{green!50!black},
    stringstyle=\color{red},
}
\usepackage{silence}
\WarningFilter{latexfont}{Font shape `T1/ntxtlf/m/up' undefined}
\WarningFilter{latexfont}{Some}
\usepackage{setspace}
\usepackage[most]{tcolorbox}
\usepackage[T1]{fontenc}
\usepackage{fourier,libertine}
\let\oldstylenums\oldstyle
\usepackage{array}
\usepackage{layout}
\usepackage{float}
\usepackage{xcolor}
\usepackage{tikz}
\usepackage{amsthm}
\usepackage{enumerate}
\usepackage{picinpar}
\pagestyle{empty}
% Tcolorboxes

\makeatother
\usepackage{thmtools}
\usepackage[framemethod=TikZ]{mdframed}
\mdfsetup{skipabove=1em,skipbelow=1em}

\newcommand{\sqed}{\hfill\includegraphics[height=2ex]{Graficas/sancoya.png}}

\theoremstyle{definition}
\declaretheoremstyle[
    headfont=\bfseries\sffamily\color{black!70!black}, bodyfont=\normalfont,
    mdframed={
        linewidth=2pt,
        rightline=false,leftline=false, topline=false, bottomline=false,
        linecolor=black, backgroundcolor=black!1!white,
    }
]{thmbox}

\declaretheoremstyle[
    headfont=\bfseries\sffamily\color{black!70!black}, bodyfont=\normalfont,
    numbered=no,
    mdframed={
        linewidth=0pt,
        rightline=false, topline=false, bottomline=false,
        linecolor=black, backgroundcolor=black!2!white,
    },
    qed=\qedsymbol
]{thmproofbox}

\declaretheoremstyle[
    headfont=\bfseries\sffamily\color{black!70!black}, bodyfont=\normalfont,
    numbered=no,
    mdframed={
        linewidth=0pt,
        rightline=false, topline=false, bottomline=false,
        linecolor=black, backgroundcolor=black!2!white,
    },
    qed=\sqed
]{sancoya}


\declaretheorem[style=thmbox, name=Definición]{definition}
\declaretheorem[style=thmbox, numbered=no, name=Ejemplo]{eg}
\declaretheorem[style=sancoya, numbered=no, name=Demostración]{sproof}
\declaretheorem[style=thmbox, name=Proposición]{prop}
\declaretheorem[style=thmbox, name=Teorema, numbered=no]{theorem}
\declaretheorem[style=thmbox, name=Lema]{lemma}
\declaretheorem[style=thmbox, numbered=no, name=Corolario]{corollary}
\declaretheorem[style=thmbox, name=Solución, numbered=no]{solution}


\declaretheorem[style=thmproofbox, name=Demostración]{replacementproof}
\renewenvironment{proof}[1][\proofname]{\vspace{-10pt}\begin{replacementproof}}{\end{replacementproof}}


\declaretheorem[style=thmbox, numbered=no, name=Nota]{note}
\declaretheorem[style=thmbox, numbered=no, ]{temp}




\usepackage{geometry}
 \geometry{
 a4paper,
 total={170mm,245mm},
 left=25mm,
 top=30mm,
 }

\usepackage{bbm}

\begin{document}

\setlength{\parindent}{0cm}
\hoffset-0.46cm
\voffset-1.46cm

\begin{window}[0,l,{\includegraphics[scale=0.21]{logo.eps}},]
\large\scshape  \hspace{0.4cm}\textsf{Universidad Nacional de Colombia} \\
\textcolor{white}{\tiny.}  \large \hspace{1.5cm} \textsf{Facultad de Ciencias} \\
\textcolor{white}{\tiny.}   \normalsize\hspace{2.2cm}\textsf{Análisis Numérico I}\\
 

\end{window}

\vspace{0.2cm}
\small
\textsf{Mateo Andrés Manosalva Amaris\\
Edgar Santiago Ochoa Quiroga\\
Sergio Alejandro Bello Torres} 
\normalsize
\dotfill
\vspace{0.7cm}

\section*{\textcolor{white}{Ejercicios}}
\vspace*{-2cm} %esto es porque me raya que la sección salga y tampoco quiero que quede espacio xd

%La madre para el que mueva algo en el main.tex
%!TEX root = ../main.tex

\section{Problema 1}
Dada la función
\[
f(x) = x + \frac{1}{x} - 2, \quad f: \mathbb{R}_{>0} \to \mathbb{R},
\]
se construye el siguiente algoritmo para aproximar la raíz $r = 1$:
\[
x_{n+1} = 2 - \frac{1}{x_n}.
\]

\begin{itemize}
    \item Verificar que si $x_0 > 1$ entonces la sucesión $\{x_n\}$ es monótona decreciente y acotada inferiormente por 1. Concluir que $x_n \to 1$, aunque esta iteración no está en las hipótesis del teorema del punto fijo. ¿Qué hipótesis no se cumple?
    \item Dar un algoritmo para aproximar la raíz de $f$ que converja cuadráticamente.
\end{itemize}

%!TEX root = ../main.tex

\subsection*{Ejercicio 2.}
Sea $\| \cdot \|$ una norma en $\mathbb{R}^n$ y $A$ una matriz invertible de tamaño $n \times n$. Pruebe que:\\

Si $Ax = b$, $(A + \delta A)(x + \delta x) = b + \delta b$ y $\|A^{-1}\| \|\delta A\| < 1$, entonces $A + \delta A$ es invertible y se cumple que:
\[
\frac{\|\delta x\|}{\|x\|} \leq \frac{\text{cond}(A)}{1 - \|A^{-1}\| \|\delta A\|} \left( \frac{\|\delta A\|}{\|A\|} + \frac{\|\delta b\|}{\|b\|} \right).
\]

\begin{sproof}
    hola x4
\end{sproof}
%!TEX root = ../main.tex

\section*{Ejercicio 3}
Sean
\[
A = \begin{bmatrix} a & 0 & 0 \\ a\delta & a & 0 \\ 0 & a\delta & a \end{bmatrix}, \quad a < 0, \delta > 0, \quad b = \begin{bmatrix} -1 \\ -1.1 \\ 0 \end{bmatrix}
\]
\begin{enumerate}
    \item[a)] Obtenga el número de condición de $A$.

    \begin{solution}
        Primero note que $A$ es invertible ya que

        $$\det A=a\det \begin{pmatrix}
            a & 0 \\
            a\delta & a \\
        \end{pmatrix}=a^3$$
    entonces como $a\neq 0$, podemos calcular el número de condición $\kappa_{\infty}(A)=\|A\|_{\infty}\|A^{-1}\|_{\infty}$, donde

    $$A^{-1}=\dfrac{1}{a}\left(
\begin{array}{ccc}
 1 & 0 & 0 \\
 -\delta & 1 & 0 \\
 \delta^2 & -\delta & 1 \\
\end{array}
\right)$$
Esto nos da que $\kappa_{\infty}(A)=a\|B\|_{\infty}\dfrac{1}{a}\|B^{-1}\|_{\infty}=\|B\|_{\infty}\|B^{-1}\|_{\infty}$, donde $B=\left(
\begin{array}{ccc}
 1 & 0 & 0 \\
 \delta  & 1 & 0 \\
 0 & \delta  & 1 \\
\end{array}
\right),$ así

$$\kappa_{\infty}(A)=\left\|\left(
\begin{array}{ccc}
 1 & 0 & 0 \\
 \delta  & 1 & 0 \\
 0 & \delta  & 1 \\
\end{array}
\right)\right\|_{\infty}\left\|\left(\begin{array}{ccc}
 1 & 0 & 0 \\
 -\delta & 1 & 0 \\
 \delta^2 & -\delta & 1 \\
\end{array}\right)\right\|_{\infty}=(1+\delta)(1+\delta+\delta^2)$$



    \end{solution}
    \item[b)] Estudie el condicionamiento del sistema $Ax = b$ en función de los valores de $\delta$. Interprete su resultado.

    \begin{solution}
        En $\delta=0$ el problema se comporta bien, hacemos este  caso ya que aunque $\delta>0$, hay que pensar en valores cercanos a $0$, en estos el problema se comporta bien, note que el conjunto solución de $(1+\delta)(1+\delta+\delta^2)>0$ es justamente $\delta>0$, y esta función es creciente, entonces conforme $\delta$ se vuelve más grande, el número de condición aumenta y nuestro problema se vuelve sencible a errores de redondeo y perturbaciones en los datos. Más precisamente

        \begin{itemize}
            \item $\lim_{\delta\to 0} \kappa_{\infty}(A)=1$.\\
            \item $\lim_{\delta  \to \infty} \kappa_\infty(A)=\infty$ 
        \end{itemize}
    \end{solution}
    \item[c)] Si $a = -1$, $\delta = 0.1$ y se considera $x^* = (1, 9/10, 1)^T$ como solución aproximada del sistema $Ax = b$ (sin obtener la solución exacta), determine un intervalo en el que esté comprendido el error relativo. ¿Es coherente con la respuesta dada en el apartado anterior?

    \begin{solution}
        Reemplazando $a=-1, \delta=0.1$ y tomando $x^* = (1, 9/10, 1)^T$ se sigue que

        $$Ax^*=\left(
\begin{array}{ccc}
 -1 & 0 & 0 \\
 -0.1 & -1 & 0 \\
 0 & -0.1 & -1 \\
\end{array}
\right)\left(
\begin{array}{c}
 1 \\
 9/10 \\
 1 \\
\end{array}
\right)=\left(
\begin{array}{c}
 -1 \\
 -1 \\
 -1.09 \\
\end{array}
\right)\approx \left(
\begin{array}{c}
 -1 \\
 -1.1 \\
 0 \\
\end{array}
\right).$$
Para determinar un intervalo del error podemos utilizar la fórmula

$$\frac{1}{\kappa_{\infty}(A)} \frac{\left\|b-b^*\right\|_{\infty}}{\|b\|_\infty} \leq \frac{\left\|x-x^*\right\|_{\infty}}{\|x\|_{\infty}} \leq \kappa_{\infty}(A) \frac{\left\|b-b^*\right\|_{\infty}}{\|b\|_{\infty}},$$
en efecto $\kappa_{\infty}(A)=(1+0.1) \left(\left(0.1\right)^2+0.1+1\right)\approx1.221$, $\|b\|_{\infty}=1.1$ y $\|b-b^*\|_{\infty}=1.09$. aplicando esto a la fórmula obtenemos

$$\frac{1}{1.211} \frac{1.09}{1.1} \leq \frac{\left\|x-x^*\right\|_{\infty}}{\|x\|_{\infty}} \leq 1.221 \frac{1.09}{1.1},$$

esto es 

\[
0.818 \leq \frac{\|x-x^*\|_{\infty}}{\|x\|_{\infty}} \leq 1.20.
\]

Esto no es incoherente con el punto anterior ya que el punto anterior me dice  que el sistema $Ax=b$ es bien condicionado para $\delta$ pequeño, en este caso $\delta=0.1$, esto no significa que al introducir un $x$ arbitrario el error deba ser pequeño, solo nos da que el sistema se comporta bien ante  perturbaciones en los datos y errores de redondeo, sin embargo si me proporcionan un $x$ erroneo como solución, esto no tiene nada que ver  con lo antes mencionado, el error no necesariamente debe ser pequeño.
    \end{solution}

    \item[d)] Si $a = -1$ y $\delta = 0.1$, ¿es convergente el método de Jacobi aplicado a la resolución del sistema $Ax = b$? Realice tres iteraciones a partir de $x_0 = (0, 0, 0)^T$.

    \begin{solution}
        Tomando $a$ y $\delta$ como fueron dados en el enunciado tenemos que 
       $$A=\begin{bmatrix}
           -1&0&0\\
           -0.1&-1&0\\
           0&-0.1&-1
       \end{bmatrix}=\begin{bmatrix}
           -1&0&0\\
           0&-1&0\\
           0&0&-1
       \end{bmatrix}-\begin{bmatrix}
           0&0&0\\
           0.1&0&0\\
           0&0.1&0
       \end{bmatrix}=D-L-U$$
       Note que tenemos que $D=-I$ y que $U=0$ en este caso, luego la matriz de iteración del método de Jacobi esta dada por
       $$T_J=D^{-1}(L+U)=\begin{bmatrix}
           -1&0&0\\
           0&-1&0\\
           0&0&-1
       \end{bmatrix}\begin{bmatrix}
            0&0&0\\
           0.1&0&0\\
           0&0.1&0
       \end{bmatrix}=\begin{bmatrix}
           0&0&0\\
           -0.1&0&0\\
           0&-0.1&0
       \end{bmatrix}$$
       Luego note que $\|T_J\|_\infty=0.1<1$,  esto nos asegura que la sucesión dada por el método converge a la solución del sistema $Ax=b.$
        Ahora para realizar las tres iteraciones del método, recordemos que la forma matricial del método de Jacobi esta dada por
        $$x^{(k+1)}=T_Jx^{(k)}+D^{-1}b$$
        Como $D^{-1}=-I$, tenemos que 
        $$D^{-1}b=\begin{bmatrix}
            1\\
            1.1\\
            0
        \end{bmatrix}$$
        Con esta información consideremos la primera iteración
        $$x^{(1)}=\begin{bmatrix}
           0&0&0\\
           -0.1&0&0\\
           0&-0.1&0
       \end{bmatrix}\begin{bmatrix}
           0\\
           0\\
           0
       \end{bmatrix}+\begin{bmatrix}
            1\\
            1.1\\
            0
        \end{bmatrix}=\begin{bmatrix}
            1\\
            1.1\\
            0
        \end{bmatrix}$$
        Ahora para la segunda iteración tenemos que
        $$x^{(2)}=\begin{bmatrix}
           0&0&0\\
           -0.1&0&0\\
           0&-0.1&0
       \end{bmatrix}\begin{bmatrix}
            1\\
            1.1\\
            0
        \end{bmatrix}+\begin{bmatrix}
            1\\
            1.1\\
            0
        \end{bmatrix}=\begin{bmatrix}
           0\\
           -0.1\\
           -0.11\\
        \end{bmatrix}+\begin{bmatrix}
            1\\
            1.1\\
            0
        \end{bmatrix}=\begin{bmatrix}
            1\\
            1\\
            -0.11
        \end{bmatrix}$$
        Para finalizar la tercera iteración esta dada por
        $$x^{(3)}=\begin{bmatrix}
           0&0&0\\
           -0.1&0&0\\
           0&-0.1&0
       \end{bmatrix}\begin{bmatrix}
            1\\
            1\\
            -0.11
        \end{bmatrix}+\begin{bmatrix}
            1\\
            1.1\\
            0
            \end{bmatrix}=\begin{bmatrix}
           0\\
           -0.1\\
           -0.1\\
        \end{bmatrix}+\begin{bmatrix}
            1\\
            1.1\\
            0
        \end{bmatrix}=\begin{bmatrix}
            1\\
            1\\
            -0.1
        \end{bmatrix}$$
        En teoría $x^{(3)}$ es una solución aproximada del sistema, pero note que si efectuamos la multiplicación de esta aproximación con $A$ tenemos que
        $$Ax^{(3)}=\begin{bmatrix}
           -1&0&0\\
           -0.1&-1&0\\
           0&-0.1&-1
       \end{bmatrix}\begin{bmatrix}
            1\\
            1\\
            -0.1
        \end{bmatrix}=\begin{bmatrix}
            -1\\
            -1.1\\
            0
        \end{bmatrix}=b.$$
        Note que justamente $x^{(3)}$ resulto ser no solo una aproximación, sino la solución exacta del sistema, es decir, que en este caso ademas de saber que el método de Jacobi converge, concluimos que convergía en tan solo 3 iteraciones.
        
    \end{solution}
\end{enumerate}
%!TEX root = ../main.tex

\subsection*{Ejercicio 4.}
Al aproximar una función continua $f : [0, 1] \to \mathbb{R}$ mediante un polinomio $p(t) = a_n t^n + \cdots + a_1 t + a_0$, el error de aproximación $E$ se mide en la norma $L^2$, es decir:
\[
E^2 := \|p - f\|_{L^2}^2 = \int_0^1 [p(t) - f(t)]^2 \, dt.
\]

\begin{enumerate}
    \item[(a)] Muestre que la minimización del error $E = E(a_0, \ldots, a_n)$ conduce a un sistema de ecuaciones lineales $H_n a = b$, donde:
    \[
    b = [b_0, \ldots, b_n]^T \in \mathbb{R}^{n+1}, \quad b_i = \int_0^1 f(t)t^i \, dt, \quad i = 0, \ldots, n,
    \]
    y $H_n$ es la matriz de Hilbert de orden $n$, definida como:
    \[
    (H_n)_{i,j} = \frac{1}{i + j + 1}, \quad i, j = 0, \ldots, n.
    \]
    El vector $a$ representa los coeficientes del polinomio $p$.\\

    \begin{proof}
    xd
    \end{proof}

    \item[(b)] Muestre que $H_n$ es simétrica y definida positiva.
    \item[(c)] Solucione el sistema $H_n x = b$, donde $b$ tiene componentes $b_i = 1 / (n + i - 1)$, para $i = 1, \ldots, n$. Para esto, use las factorizaciones LU ($[L, U] = \text{lu}(H)$) y Cholesky ($L = \text{chol}(H)$). Luego resuelva los dos sistemas triangulares.
    \item[(d)] Para ambos métodos, ¿qué tan precisas son las soluciones numéricas $\hat{x}_{\text{approx}}$? Tabule los errores de la solución:
    \[
    e(n) = \|\hat{x}_{\text{approx}} - x_{\text{exact}}\|
    \]
    como una función de $n = 2, \ldots, 15$. Note que $x_{\text{exact}} = (0, \ldots, 1)^T$. Puede graficar los errores en función de $n$ utilizando la función \texttt{semilogy} de Matlab. Explique en detalle los resultados.
\end{enumerate}


%!TEX root = ../main.tex

\item Si $f: X_1 \rightarrow X_2$ es un homeomorfismo entre espacios Hausdorff localmente compactos, demuestra que $f$ se extiende a un homeomorfismo de sus compactificaciones por un punto.

\begin{proof}
    Suponga $f:X_1\to X_2$ un homeomorfismo y definamos

    \begin{align*}
        g : Y_1 &\longrightarrow Y_2 \\
        x &\longmapsto g(x) = \begin{cases}
            f(x) &\text{si } x\in X_1\\
            \tilde{y} &\text{si }x=y,
        \end{cases}
    \end{align*}

    con $Y_1, Y_2$ las compactificaciones a un punto, $g$ es un homeomorfismo entre $Y_1$ y $Y_2$. Sea $V$ un abierto, $V\subset Y_2$. Tenemos dos casos, $V$ es abierto de $X_2$ o $V=Y_2-C$ con $C$ compacto en $X_2$, el primer caso es trivial. Para el segundo caso note que 

    \begin{align*}
        g^{-1}(V)&=g^{-1}(Y_2-C)\\
        &=Y_1-g^{-1}(C)\\
        &=Y_1- K_1
    .\end{align*}

    con $K_1$ compacto en $X_1$. Análogamente 

    $$g(Y_1-C)=Y_2-g(C)=Y_2-K_2$$

    con $K_2$ compacto en $X_2$.
        
\end{proof}


%!TEX root = ../main.tex

\item Demuestra que la compactificación por un punto de $\mathbb{R}$ es homeomorfa al círculo $S^1$.\\

\begin{proof}
    Vamos a construir el homeomorfismo, este es la proyección estereográfica.

\begin{center}
    \tikzset{every picture/.style={line width=1.3pt}}

\begin{tikzpicture}[x=2pt,y=2pt,yscale=-1,xscale=1]
%uncomment if require: \path (0,300); %set diagram left start at 0, and has height of 300

%Shape: Ellipse [id:dp2076532777529081] 
\draw   (141.83,89.08) .. controls (141.83,68.14) and (158.81,51.17) .. (179.75,51.17) .. controls (200.69,51.17) and (217.67,68.14) .. (217.67,89.08) .. controls (217.67,110.02) and (200.69,127) .. (179.75,127) .. controls (158.81,127) and (141.83,110.02) .. (141.83,89.08) -- cycle ;
%Shape: Ellipse [id:dp49047118671320233] 
\draw  [draw opacity=0][fill={rgb, 255:red, 208; green, 2; blue, 27 }  ,fill opacity=1 ] (178.97,89.08) .. controls (178.97,88.65) and (179.32,88.3) .. (179.75,88.3) .. controls (180.18,88.3) and (180.53,88.65) .. (180.53,89.08) .. controls (180.53,89.52) and (180.18,89.87) .. (179.75,89.87) .. controls (179.32,89.87) and (178.97,89.52) .. (178.97,89.08) -- cycle ;
%Straight Lines [id:da9873077654790294] 
\draw    (179.73,31.6) -- (180.33,143.67) ;
%Straight Lines [id:da48767165474320806] 
\draw    (120.8,88.93) -- (263.13,88.93) ;
%Straight Lines [id:da6624476477384516] 
\draw    (179.75,51.17) -- (250.13,88.93) ;
%Shape: Circle [id:dp8959989266574702] 
\draw  [draw opacity=0][fill={rgb, 255:red, 208; green, 2; blue, 27 }  ,fill opacity=1 ] (210.6,68.43) .. controls (210.6,67.79) and (211.12,67.27) .. (211.77,67.27) .. controls (212.41,67.27) and (212.93,67.79) .. (212.93,68.43) .. controls (212.93,69.08) and (212.41,69.6) .. (211.77,69.6) .. controls (211.12,69.6) and (210.6,69.08) .. (210.6,68.43) -- cycle ;

% Text Node
\draw (181.2,43.07) node [anchor=north west][inner sep=0.75pt]  [font=\tiny,xscale=2,yscale=2]  {$1$};
% Text Node
\draw (247.33,91.82) node [anchor=north west][inner sep=0.75pt]  [font=\tiny,xscale=2,yscale=2]  {$a$};


\end{tikzpicture}
\end{center}

<<<<<<< HEAD
=======
Por la ecuación de punto pendiente, la recta que se observa en el dibujo y que corta la circunferencia  tiene la ecuación 

$$y=\frac{-x}{a}+1,$$

reemplazando esto en la ecuación del círculo $(x^2+y^2=1)$obtenemos que 

$$x^2+\frac{x^2}{a^2}-\frac{2x}{a}=0,$$

así pues

$$x\left(1+\frac{1}{a^2}\right)=\frac{2}{a},$$

esto es 

$$x=\frac{2a}{a^2+1}.$$

Finalmente, al reemplazar $x$ en la ecuación de la recta, obtenemos que 

$$y=\frac{-2}{a^2+1}+1$$

que es el heomemorfimo ya que cuando $a\to \infty$ obtenemos el polo norte y la función establece una correspondencia entre un punto en la circunferencia y un punto en la recta (por la construcción que hicimos). Obviamente la función es biyectiva y continua y su inversa es continua xd.

\end{proof}
>>>>>>> fa5901a0b37caa9a02deddf9cef5b61772a49d29


\end{document}